\documentclass[12pt]{article}
\usepackage[utf8]{inputenc}
\usepackage[margin=0.2in]{geometry}
\usepackage{tcolorbox}
\usepackage{amsmath}
\usepackage{ulem}
\usepackage{enumitem}
\usepackage{amssymb}

\tcbset{
    frame code={},
    center title,
    left=0pt,
    right=0pt,
    top=0pt,
    bottom=0pt,
    colback=white!10,
    colframe=white!20,
    width=\textwidth,
    enlarge left by=0mm,
    boxsep=3pt,
    arc=0pt,
    outer arc=0pt,
    before skip=0pt,
    after skip=0pt,
}

\setlength{\parindent}{0pt}
\setlength{\parskip}{0pt}
\pagestyle{empty}

\begin{document}

\begin{center}
\large\textbf{CSC112 - TEST2: 8-Bit Binary Arithmetic} (\textbf{Marks}: 25) \\
 \textbf{ID}: \underline{\hspace{3cm}} \textbf{Name}: \underline{\hspace{8cm}} \quad  \textbf{Date}: 2025-12-08  
\end{center}

\textbf{Instructions:} For all questions, assume an 8-bit binary machine where the $8^{th}$ bit is the sign bit (0 for positive, 1 for negative). Only 7 bits are available for magnitude. Given: $A = 20$, $B = 15$, $C = -9$, $D = -12$, answer the following:-

\begin{tcolorbox}
(a) Calculate binary sum of A and B. \hfill  \uline{\textbf{ANSWER}: \hspace{0.2\textwidth}}

\vspace{4cm}

\end{tcolorbox}
\begin{tcolorbox}
(b) Calculate A - B using 1's complement method.  \hfill  \uline{\textbf{ANSWER}: \hspace{0.2\textwidth}}

\vspace{4cm}

\end{tcolorbox}
\begin{tcolorbox}
(c) Calculate A - B using 2's complement method. \hfill  \uline{\textbf{ANSWER}: \hspace{0.2\textwidth}}

\vspace{4cm}

\end{tcolorbox}
\begin{tcolorbox}
(d) Compute C + D using 2's complement method. \hfill  \uline{\textbf{ANSWER}: \hspace{0.2\textwidth}}

\vspace{4cm}

\end{tcolorbox}
\begin{tcolorbox}
(e) Can $130_{10}$ be represented in the 8-bit machine? Justify. \hfill \uline{\textbf{ANSWER}: \hspace{0.2\textwidth}}

\vspace{4cm}

\end{tcolorbox}

\end{document}